%! Rates of Change — Unit 1, Lesson 1.3 — Homework (Answer Key)
\documentclass[10pt]{article}
\usepackage{precalculus-article}
\usepackage{precalculus-key}   % pulls in -boxes; do NOT also load -boxes
\newcommand{\TallMath}[1]{$\displaystyle #1\rule[-1.4em]{0pt}{3.2em}$}

\begin{document}

\pageheader{Unit 1, Lesson 1.3}{Homework --- Answer Key}

\vspace{0.12in}

\begin{enumerate}[itemsep=10pt]

\item Let $g(x) = x^2 - 4x$. Find the average rate of change over each
      interval.

\begin{enumerate}[label=(\alph*), itemsep=6pt]
  \item Over $[1, 4]$:

        \begin{work}
          g(1) &= (1)^2 - 4(1) \\
               &= -3 \\
          g(4) &= (4)^2 - 4(4) \\
               &= 0 \\
          \text{AROC} &= \frac{0 - (-3)}{4 - 1} \\
                      &= 1
        \end{work}

        AROC $= $ \ans{1}

  \item Over $[0, 2]$:

        \begin{work}
          g(0) &= 0 \\
          g(2) &= (2)^2 - 4(2) \\
               &= -4 \\
          \text{AROC} &= \frac{-4 - 0}{2 - 0} \\
                      &= -2
        \end{work}

        AROC $= $ \ans{$-2$}
\end{enumerate}

\item The table gives a town's population $P(t)$.

\smallskip
\begin{tabular}{|c|c|c|c|c|}
\hline
Year & 1990 & 2000 & 2010 & 2020 \\
\hline
$P(t)$ & 12{,}000 & 15{,}000 & 21{,}000 & 23{,}000 \\
\hline
\end{tabular}
\smallskip

\begin{enumerate}[label=(\alph*), itemsep=4pt]
  \item AROC from 1990 to 2000: \ans{300} people per year
  \item AROC from 2010 to 2020: \ans{200} people per year
  \item Which decade had the fastest growth? \quad \ans{2000--2010 (600/yr)}
\end{enumerate}

\item A hiker's elevation $E(t)$, in feet $t$ hours into a hike, has an
      average rate of change of $-40$ feet per hour over $[2, 5]$. Write one
      sentence explaining what this means about the hike.

      \vspace{0.05in}
      \ansline{Between hour 2 and hour 5 the hiker descended, losing an average of 40 feet}
      \ansline{of elevation each hour --- negative means elevation drops as time increases.}

\boxguard[16]
\item A function $q$ has the table below. Fill in the AROC row, then answer.

\smallskip
\renewcommand{\arraystretch}{1.5}
\begin{tabular}{|c|c|c|c|c|c|}
\hline
$x$ & 0 & 1 & 2 & 3 & 4 \\
\hline
$q(x)$ & 1 & 2 & 5 & 10 & 17 \\
\hline
\end{tabular}
\smallskip

AROCs over $[0,1], [1,2], [2,3], [3,4]$: \quad
\ans{1} , \ans{3} , \ans{5} , \ans{7}

\begin{enumerate}[label=(\alph*), itemsep=4pt]
  \item Each AROC changes by \ans{2} --- so $q$ is
        \textbf{linear} / \ans{quadratic} (circle one).
  \item The AROCs are increasing, so the graph of $q$ is concave
        \ans{up}.
\end{enumerate}

\item The table shows values of a function $f$. What is the average rate of
      change of $f$ over $[3, 7]$?

\smallskip
\begin{tabular}{|c|c|c|c|c|}
\hline
$x$ & 1 & 3 & 5 & 7 \\
\hline
$f(x)$ & 4 & 10 & 22 & 40 \\
\hline
\end{tabular}
\smallskip

\begin{enumerate}[label=\Alph*., itemsep=3pt]
  \item $7.5$ \quad \textcolor{keyred}{\textbf{$\leftarrow$ correct: $(40-10)/(7-3)$}}
  \item $30$
  \item $6$
  \item $10$
\end{enumerate}

\end{enumerate}

\vspace{0.1in}

\boxguard
\begin{extensionbox}
For $f(x) = x^2$, use $f(3.1) = 9.61$ and $f(3.01) = 9.0601$ to compute the
AROC over $[3, 3.1]$ and over $[3, 3.01]$. The answers creep toward one
value --- the rate of change \textit{at} $x = 3$. What is it?

\vspace{0.05in}
\ansline{AROC over $[3, 3.1]$: $0.61/0.1 = 6.1$. AROC over $[3, 3.01]$: $0.0601/0.01 = 6.01$.}
\ansline{The values creep toward 6 --- the rate of change at $x = 3$.}
\end{extensionbox}

\vspace{0.1in}

\boxguard
\begin{notesbox}{Preview of Lesson 1.4}
Rates measure how a graph moves; next we ask how a whole graph can be
\textit{moved}. Lesson 1.4, \textit{Transformations of Functions}, slides,
flips, and stretches graphs --- and shows how the formula records each move.
\end{notesbox}

\end{document}
